\section{Demostraciones complementarias de SP2}
\label{app:sp2-proofs}

La alineación marginal de SP2 se demuestra para preferencias continuas, umbrales positivos y una matriz de servicio fija durante cada decisión. El resultado identifica el gradiente del potencial con el pago marginal. El cierre entero, la mecánica y la ejecución vecinal discretizada se tratan mediante resultados independientes.

\subsection{Potencial de la puntuación marginal y frontera de integrabilidad}
\label{app:sp2-marginal-proof}

La regla de la cadena aplicada al potencial de~\eqref{eq:sp2-potential} da
\[
  \frac{\partial\Phi}{\partial\rho_{ik}}
  =\frac{e_{ik}}{d_k^{\mathrm{srv}}}V_k
   \left[1-\frac{S_k(\rho)}{d_k^{\mathrm{srv}}}\right]_+-g_{ik}
  =f_{ik}^{\mathrm{marg}}.
\]
La puntuación marginal coincide componente a componente con el gradiente de $\Phi$ en cada región diferenciable.

En el interior no saturado, las derivadas cruzadas de la puntuación plana son
\[
  \frac{\partial f_{ik}^{\mathrm{plain}}}{\partial\rho_{jk}}
  =-\frac{V_ke_{jk}}{d_k^{\mathrm{srv}}},
  \qquad
  \frac{\partial f_{jk}^{\mathrm{plain}}}{\partial\rho_{ik}}
  =-\frac{V_ke_{ik}}{d_k^{\mathrm{srv}}}.
\]
Si $e_{ik}\neq e_{jk}$, estas expresiones difieren y violan la condición de simetría necesaria para que el campo plano sea el gradiente de un potencial escalar dos veces diferenciable en esa región. Para la puntuación marginal, en cambio, ambas derivadas cruzadas valen
\[
  -\frac{V_ke_{ik}e_{jk}}{(d_k^{\mathrm{srv}})^2}.
\]

La función $[1-s]_+$ es continua en $s=1$. Su integral también es continua y diferenciable, aunque su segunda derivada cambie por tramos. El potencial de~\eqref{eq:sp2-potential} enlaza de forma continua las regiones $S_k<d_k^{\mathrm{srv}}$ y $S_k\geq d_k^{\mathrm{srv}}$. En ambas regiones, y mediante la derivada correspondiente en la frontera, su gradiente coincide con la puntuación marginal.
